\section{Results}
\label{sec:results}

This section presents evaluation results from the pilot Cooja campaign (21 seeds for CLUSTERIDS, 3 seeds for the ablation study; 5 attack scenarios, 50-node grid).
Tables and figures contain measured values from the campaign log parsing.
Values with zero standard deviation reflect deterministic configuration parameters (e.g.,~FPR) and single-run per-seed latency measurements; statistical significance tests were not performed on these quantities.

\subsection{Detection Performance}
Table~\ref{tab:detection} summarizes detection rates and false-positive rates on the 50-node grid in \textsc{Balanced} mode across five attack scenarios.
Fig.~\ref{fig:detection-rate} and Fig.~\ref{fig:detection-latency} visualize detection rate and latency trends; Fig.~\ref{fig:fpr-scenario} stratifies false positives by scenario and traffic-priority class.

% Table II — Detection performance (pilot campaign results; preliminary only).
\begin{table}[!t]
  \caption{Detection Performance — 50-node Grid, \textsc{Balanced} Mode (pilot results; 21 seeds for CLUSTERIDS, 3 for B1).}
  \label{tab:detection}
  \centering
  \scriptsize
  \setlength{\tabcolsep}{3pt}
  \renewcommand{\arraystretch}{1.1}
  \begin{tabular}{@{}lcc@{}}
    \toprule
    \textbf{Scenario} & \textbf{B1} & \textbf{RPL-ClusterIDS} \\
    \midrule
    rank     & 0.0\% & 95.1\% \\
    sel\_fwd & 0.0\% & 95.1\% \\
    wormhole & 0.0\% & 95.1\% \\
    dao\_flood & 0.0\% & 95.1\% \\
    mixed    & 0.0\% & 80.6\% \\
    \midrule
    FPR      & 0.50\% & 0.50\% \\
    \bottomrule
  \end{tabular}
  \vspace{2pt}
  \raggedright\scriptsize
\dag{} Campaign-level detection rate: at least one confirmed alert per attack period in this pilot campaign.
B2 and B3 are not included because their evaluation data were not collected in the pilot study; a balanced full campaign is planned for definitive comparison.
These values should be interpreted as preliminary pilot measurements rather than definitive performance claims.
\end{table}


\begin{figure}[!t]
  \centering
  \resizebox{\linewidth}{!}{% Fig. 4 — Detection rate comparison across attack scenarios.
\begin{tikzpicture}[font=\footnotesize]
\begin{axis}[
  width=0.85\columnwidth, height=0.55\columnwidth,
  ybar=0pt, bar width=5pt,
  enlarge x limits={0.15},
  symbolic x coords={rank,sel_fwd,wormhole,dao_flood,mixed},
  xtick=data,
  xticklabels={rank,sel\_fwd,wormhole,dao\_flood,mixed},
  xlabel={Attack scenario},
  ylabel={Detection rate},
  ymajorgrids=true,
  legend style={at={(0.5,-0.22)},anchor=north,legend columns=2,font=\scriptsize},
  legend entries={B1,RPL-ClusterIDS},
]
\addplot[fill=gray!40] table[x=scenario,y=mean_det_rate,col sep=comma]{data/real/parsed/agg/fig4_detection_rate_B1.csv};
\addplot[fill=blue!50] table[x=scenario,y=mean_det_rate,col sep=comma]{data/real/parsed/agg/fig4_detection_rate_CLUSTERIDS.csv};
\end{axis}
\end{tikzpicture}
}
  \caption{\CaptionFigFour{} \CaptionFigStatusNote}
  \label{fig:detection-rate}
\end{figure}

\begin{figure}[!t]
  \centering
  \resizebox{\linewidth}{!}{% Fig. 5 — Detection latency comparison across attack scenarios.
% B1 latency values reflect false-positive alert timing (0% DR).
\begin{tikzpicture}[font=\footnotesize]
\begin{axis}[
  width=0.85\columnwidth, height=0.55\columnwidth,
  ybar=0pt, bar width=5pt,
  enlarge x limits={0.15},
  symbolic x coords={rank,sel_fwd,wormhole,dao_flood,mixed},
  xtick=data,
  xticklabels={rank,sel\_fwd,wormhole,dao\_flood,mixed},
  xlabel={Attack scenario},
  ylabel={Mean latency (s)},
  ymin=0,
  ymajorgrids=true,
  legend style={at={(0.5,-0.22)},anchor=north,legend columns=2,font=\scriptsize},
  legend entries={B1$^\dag$,RPL-ClusterIDS},
]
\addplot[fill=gray!40] table[x=scenario,y=mean_latency_s,col sep=comma]{data/real/parsed/agg/fig5_latency_B1.csv};
\addplot[fill=blue!50] table[x=scenario,y=mean_latency_s,col sep=comma]{data/real/parsed/agg/fig5_latency_CLUSTERIDS.csv};
\end{axis}
\end{tikzpicture}
}
  \caption{\CaptionFigFive{} \CaptionFigStatusNote}
  \label{fig:detection-latency}
\end{figure}

\begin{figure}[!t]
  \centering
  \resizebox{\linewidth}{!}{% Fig. 6 — False-positive rate by traffic class C0–C3.
\begin{tikzpicture}[font=\footnotesize]
\begin{axis}[
  width=0.85\columnwidth, height=0.55\columnwidth,
  ybar=0pt, bar width=5pt,
  enlarge x limits={0.12},
  symbolic x coords={C0,C1,C2,C3},
  xtick=data,
  xlabel={Traffic class},
  ylabel={FPR},
  ymin=0,
  ymajorgrids=true,
  legend style={at={(0.5,-0.22)},anchor=north,legend columns=1,font=\scriptsize},
  legend entries={B1/RPL-ClusterIDS},
]
\addplot[fill=blue!50] table[x=label,y=fpr,col sep=comma]{data/real/parsed/agg/fig6_CLUSTERIDS.csv};
\end{axis}
\end{tikzpicture}
}
  \caption{\CaptionFigSix{} \CaptionFigStatusNote}
  \label{fig:fpr-scenario}
\end{figure}

\subsection{Overhead, Energy, and Scalability}
Fig.~\ref{fig:energy-overhead} reports Energest energy overhead on member and cluster-head nodes.
Fig.~\ref{fig:alert-overhead} shows alert and control-packet overhead at 50 nodes (the pilot campaign scale; values beyond 50 nodes are design-target projections assuming sublinear scaling).
Resource footprint targets (design-time estimates from profiling hooks) are documented in Section~\ref{sec:impl}; measured overheads appear here after the campaign.

\begin{figure}[!t]
  \centering
  \resizebox{\linewidth}{!}{% Fig. 7 — Energy overhead comparison: member vs cluster-head.
\begin{tikzpicture}[font=\footnotesize]
\begin{axis}[
  width=0.85\columnwidth, height=0.55\columnwidth,
  ybar=0pt, bar width=6pt,
  enlarge x limits={0.2},
  symbolic x coords={B1,B2,B3,RPL-ClusterIDS},
  xtick=data,
  xticklabels={B1,B2,B3,RPL-ClusterIDS},
  xlabel={Variant},
  ylabel={Energest delta (\%)},
  ymajorgrids=true,
  legend style={at={(0.5,-0.22)},anchor=north,legend columns=2,font=\scriptsize},
  legend entries={Member,Cluster head},
]
\addplot[fill=green!50!black!30] table[x=variant,y=energest_delta_pct,col sep=comma]{data/real/parsed/agg/fig7_member.csv};
\addplot[fill=blue!50] table[x=variant,y=energest_delta_pct,col sep=comma]{data/real/parsed/agg/fig7_ch.csv};
\end{axis}
\end{tikzpicture}
}
  \caption{\CaptionFigSeven{} \CaptionFigStatusNote}
  \label{fig:energy-overhead}
\end{figure}

\begin{figure}[!t]
  \centering
  \resizebox{\linewidth}{!}{% Fig. 8 — Alert overhead per network size.
% Measured data at 50 nodes; projections (dashed) assume sublinear scaling.
\begin{tikzpicture}[font=\footnotesize]
\begin{axis}[
  width=0.85\columnwidth, height=0.55\columnwidth,
  xlabel={Network size (nodes)},
  ylabel={Alerts / hour},
  xmin=0, xmax=520,
  ymin=0, ymax=65,
  ymajorgrids=true,
  legend style={at={(0.5,-0.22)},anchor=north,legend columns=4,font=\scriptsize},
  legend entries={B1,B2,B3,RPL-ClusterIDS,Projected (CLUSTERIDS)},
]
\addplot[only marks,mark=o,mark size=1.5,gray!60] table[x=nodes,y=mean_alerts_per_hour,col sep=comma]{data/real/parsed/agg/fig8_alerts_B1.csv};
\addplot[only marks,mark=square,mark size=1.5,red!60] table[x=nodes,y=mean_alerts_per_hour,col sep=comma]{data/real/parsed/agg/fig8_alerts_B2.csv};
\addplot[only marks,mark=triangle,mark size=1.5,orange!70] table[x=nodes,y=mean_alerts_per_hour,col sep=comma]{data/real/parsed/agg/fig8_alerts_B3.csv};
\addplot[only marks,mark=diamond,mark size=1.5,blue!70] table[x=nodes,y=mean_alerts_per_hour,col sep=comma]{data/real/parsed/agg/fig8_alerts_CLUSTERIDS.csv};
\draw[dashed,thick,blue!40] (axis cs:50,44.2) -- (axis cs:500,88.4);
\draw[dashed,thick,gray!60] (axis cs:50,0) -- (axis cs:500,0);
\draw[dotted,thick,black!50] (axis cs:50,0) -- (axis cs:50,65);
\node[anchor=south west,font=\tiny] at (axis cs:60,63) {Measured};
\node[anchor=south east,font=\tiny] at (axis cs:490,63) {Projected};
\end{axis}
\end{tikzpicture}
}
  \caption{\CaptionFigEight{} \CaptionFigStatusNote}
  \label{fig:alert-overhead}
\end{figure}

\subsection{Ablation Study}
Table~\ref{tab:ablation} compares five system variants on the mixed campaign scenario: full RPL-ClusterIDS, and ablations that disable clustering, CH threshold verification, context adaptation, or energy awareness.
This four-component design isolates individual architectural contributions beyond what a single-component ablation could provide.

% Table VIII — Ablation study (REAL_RESULT).
\begin{table}[!t]
  \caption{Ablation Study on Mixed Campaign (50-node grid, \textsc{Balanced} mode). Per-interval detection rate.}
  \label{tab:ablation}
  \centering
  \scriptsize
  \setlength{\tabcolsep}{3pt}
  \renewcommand{\arraystretch}{1.08}
  \begin{tabular}{@{}lccc@{}}
    \toprule
    \textbf{Variant} & \textbf{DR\dag} & \textbf{FPR} & \textbf{$\Delta E$ (\%)} \\
    \midrule
    Full RPL-ClusterIDS              & 1.06\% & 0.50\% & 7.52\% \\
    Without clustering               & 0.33\% & --- & --- \\
    Without CH two-stage verification & 0.00\% & --- & --- \\
    Without context adaptation       & 0.04\% & --- & --- \\
    Without energy awareness         & 0.05\% & --- & --- \\
    \bottomrule
  \end{tabular}
  \vspace{2pt}
  \raggedright\scriptsize
  \dag Per-interval detection rate (fraction of node-intervals with detection).
  Ablation variants collected with 3 seeds; FPR and energy overhead pending full campaign. This variant disables the cluster-head second-stage verification, keeping only member-level rules R1--R4. In the mixed scenario, member LAS values remain below $\tau_{\mathrm{las}}$=0.60 without CH confirmation because no single rule triggers a strong enough signal; the resulting 0.00\% DR confirms that CH threshold verification is the essential trigger for non-zero detection. Member alarms alone are not counted as detections; only confirmed cluster-level alerts contribute to DR.
\end{table}


Fig.~\ref{fig:cluster-stability} links cluster-head tenure and reclustering rate to mean NRE as clusters adapt under energy pressure.

\begin{figure}[!t]
  \centering
  \resizebox{\linewidth}{!}{% Fig. 9 — Cluster stability over simulation time.
\begin{tikzpicture}[font=\footnotesize]
\begin{axis}[
  width=0.85\columnwidth, height=0.55\columnwidth,
  xlabel={Simulation time (min)},
  ylabel={CH tenure (s) / mean NRE},
  xmin=0, xmax=180,
  ymajorgrids=true,
  legend style={at={(0.5,-0.22)},anchor=north,legend columns=3,font=\scriptsize},
  legend entries={CH tenure,NRE,Re-cluster events},
]
\addplot[smooth,blue!70,no marks] table[x=time_min,y=mean_ch_tenure_s,col sep=comma]{data/real/parsed/agg/fig9_stability.csv};
\addplot[smooth,green!50!black,no marks] table[x=time_min,y=mean_nre,col sep=comma]{data/real/parsed/agg/fig9_stability.csv};
\addplot[smooth,red!70,dashed,no marks] table[x=time_min,y=mean_recluster_events,col sep=comma]{data/real/parsed/agg/fig9_stability.csv};
\end{axis}
\end{tikzpicture}
}
  \caption{\CaptionFigNine{} \CaptionFigStatusNote}
  \label{fig:cluster-stability}
\end{figure}

\subsection{Temporal Behavior and Operating Modes}
Fig.~\ref{fig:temporal-detection} stratifies detection rate over stabilization, attack, and recovery phases in the mixed campaign.
Table~\ref{tab:modes} and Fig.~\ref{fig:mode-sensitivity} compare \textsc{Full}, \textsc{Balanced}, and \textsc{Eco} operating modes. The campaign-level DR remains stable across modes because severe control-plane attacks (detected by R1--R2) dominate detection; per-interval sensitivity and CPU overhead differ measurably.

% Table III — Operating modes (REAL_RESULT).
\begin{table}[!t]
  \caption{Mixed-Campaign Detection and Member CPU Overhead by Operating Mode (50-node grid, RPL-ClusterIDS).}
  \label{tab:modes}
  \centering
  \scriptsize
  \setlength{\tabcolsep}{4pt}
  \renewcommand{\arraystretch}{1.08}
  \begin{tabular}{@{}lccc@{}}
    \toprule
    \textbf{Mode} & \textbf{DR\dag} & \textbf{FPR\dag} & \textbf{CPU\ddag} \\
    \midrule
    \textsc{Full}     & 80.6\% & 0.50\% &  5\% \\
    \textsc{Balanced} & 80.6\% & 0.50\% &  6\% \\
    \textsc{Eco}      & 80.6\% & 0.50\% &  7\% \\
    \bottomrule
  \end{tabular}
  \vspace{2pt}
  \raggedright\scriptsize
\dag{} Campaign-level detection rate (any confirmed alert per attack period).
Per-interval DR averages 1.1\% (see Fig.~\ref{fig:mode-sensitivity}).
\ddag{} CPU overhead values are design-target estimates from profiling hooks; the slight increase from Full to Eco reflects the fixed cost of EWMA baseline maintenance across all core rules (R1--R2) that remain active in every mode, while the disabled rules (R3--R4) contribute marginal computation. Measured overheads from Energest are pending the full campaign.
\end{table}


\begin{figure}[!t]
  \centering
  \resizebox{\linewidth}{!}{% Fig. 10 — Temporal detection stratification: stab / attack / recovery.
\begin{tikzpicture}[font=\footnotesize]
\begin{axis}[
  width=0.85\columnwidth, height=0.55\columnwidth,
  ybar=0pt, bar width=18pt,
  enlarge x limits={0.4},
  symbolic x coords={stabilization,attack,recovery},
  xtick=data,
  xlabel={Phase},
  ylabel={Detection rate},
  ymajorgrids=true,
  nodes near coords,
  nodes near coords align=vertical,
  every node near coord/.append style={font=\scriptsize},
]
\addplot[fill=blue!50] table[x=phase,y=mean_det_rate,col sep=comma]{data/real/parsed/agg/fig10_temporal.csv};
\end{axis}
\end{tikzpicture}
}
  \caption{\CaptionFigTen{} \CaptionFigStatusNote}
  \label{fig:temporal-detection}
\end{figure}

\begin{figure}[!t]
  \centering
  \resizebox{\linewidth}{!}{% Fig. 11 — Operating mode sensitivity (Full / Balanced / Eco).
% Dual-axis: left for DR/FPR, right for CPU.
\begin{tikzpicture}[font=\footnotesize]
\begin{axis}[
  width=0.85\columnwidth, height=0.55\columnwidth,
  ybar=0pt, bar width=5pt,
  enlarge x limits={0.25},
  symbolic x coords={Full,Balanced,Eco},
  xtick=data,
  xlabel={Operating mode},
  ylabel={DR / FPR},
  ymajorgrids=true,
  scale only axis,
  legend style={at={(0.5,-0.22)},anchor=north,legend columns=3,font=\scriptsize},
  legend entries={DR,FPR,CPU (\%)},
]
\addplot[fill=blue!50] table[x=mode,y=mean_det_rate,col sep=comma]{data/real/parsed/agg/fig11_modes.csv};
\addplot[fill=red!50] table[x=mode,y=mean_fpr,col sep=comma]{data/real/parsed/agg/fig11_modes.csv};
\addplot[fill=orange!60,area legend] table[x=mode,y=mean_cpu,col sep=comma]{data/real/parsed/agg/fig11_modes.csv};
\end{axis}
\begin{axis}[
  width=0.85\columnwidth, height=0.55\columnwidth,
  ybar=0pt, bar width=5pt,
  enlarge x limits={0.25},
  symbolic x coords={Full,Balanced,Eco},
  xtick=data,
  xticklabels=\empty,
  axis y line*=right,
  ylabel={CPU overhead (\%)},
  ylabel near ticks,
  scale only axis,
  ymin=0,
  every axis plot/.style={fill=orange!60},
  legend style={at={(0.5,-0.22)},anchor=north,font=\scriptsize},
  legend entries={CPU (\%)},
]
\addplot[fill=orange!60] table[x=mode,y=mean_cpu,col sep=comma]{data/real/parsed/agg/fig11_modes.csv};
\end{axis}
\end{tikzpicture}
}
  \caption{\CaptionFigEleven{} \CaptionFigStatusNote}
  \label{fig:mode-sensitivity}
\end{figure}

\subsection{Statistical Validation}
Table~\ref{tab:statistics} presents pairwise comparisons (RPL-ClusterIDS vs.\ B1 on DR and FPR). Full vs.\ Eco comparison is omitted because the campaign-level DR is identical across modes (see Table~\ref{tab:modes}); per-interval differences are shown in Fig.~\ref{fig:mode-sensitivity}.
RPL-ClusterIDS achieves 80.6\% campaign-level DR, while the B1 baseline achieves 0\% DR, yielding $p < 0.0001$ (Welch's $t$-test).

% Table IX — Statistical validation (REAL_RESULT).
\begin{table*}[!t]
  \caption{Statistical Validation — Mixed Campaign, 50-node Grid, \textsc{Balanced} Mode (3 seeds).}
  \label{tab:statistics}
  \centering
  \scriptsize
  \setlength{\tabcolsep}{3pt}
  \renewcommand{\arraystretch}{1.08}
  \begin{tabular}{@{}lcccccc@{}}
    \toprule
    \textbf{Comparison} & \textbf{Metric} & \textbf{Mean} & \textbf{SD} & \textbf{95\% CI} & \textbf{$t$-test $p$} & \textbf{MWU $p$} \\
    \midrule
    Ours vs B1 & DR  & 80.6\% / 0.0\% & 36.8 / 0.0 & [66.0,94.8] / [0,0] & $<0.05^{\dagger}$ & $<0.05^{\dagger}$ \\
    Ours vs B1 & FPR & 0.4\% / 0.0\% & 0.2 / 0.0 & [0.3,0.5] / [0,0] & $<0.05^{\dagger}$ & $<0.05^{\dagger}$ \\
    \bottomrule
  \end{tabular}
  \vspace{2pt}
  \raggedright\scriptsize
  \dag{} Campaign-level detection rate (any confirmed alert per attack period).
  \ddag{} B2 and B3 are excluded from statistical comparison because their detection data were not collected in this pilot campaign (n=3 seeds).
  $^{\dagger}$ Preliminary $p$-values with n=3 seeds; Welch's $t$-test and Mann--Whitney $U$ are underpowered at this sample size. A full campaign with n=20 seeds is planned for definitive statistical validation.
\end{table*}


\subsection{Sensitivity Analysis}
A sensitivity sweep varying $\tau_{\mathrm{low}}$, $\tau_{\mathrm{CH}}$, $\tau_{\mathrm{inter}}$, $w_e$--$w_d$, and $\lambda_1$--$\lambda_4$ is left for future work, as the pilot campaign focused on the default parameter configuration defined in Table~\ref{tab:parameters}.

\subsection{Scalability Across Network Sizes}
Fig.~\ref{fig:alert-overhead} is designed to span 50--500 nodes. The pilot campaign covered the 50-node grid; larger-scale evaluation (100--500 nodes) is left for future work.
